\section{Open questions}

The following five questions are genuinely open after this replication
and are the natural next moves. All can be attacked with free-tier
compute (local Qiskit + Argo).

\begin{enumerate}

\item \textbf{Time-dependent Hamiltonians.}
Does the Chebyshev-interpolation advantage survive on time-dependent
$H(t)$, where the effective Trotter Hamiltonian is no longer even in $s$
and the reflection-symmetry trick cannot be used?
\emph{Basis.} Paper only proves smoothness/evenness of $\tilde H_s$ at
$s=0$ for time-independent $H$ (Sec 2--3). For $H(t)$ the symmetric-ST
evenness argument breaks; a modified interpolation scheme (e.g.,
asymmetric nodes, or Magnus-expansion-based $\tilde H_s$) may or may
not still give spectral convergence.
\emph{Next steps.} Take a driven TFIM
$H(t) = -J(ZZ + g(t)(XI+IX))$ with $g(t) = g_0 + A\sin(\omega t)$, and
compare interpolation of Magnus-truncated $\tilde H_s$ vs unsymmetrized
Chebyshev on raw fractional queries. Look at slope on error-vs-$n$;
quantify order lost if slope collapses to polynomial. Purely CPU / free.

\item \textbf{Composition with LCU / qubitization.}
Can Chebyshev-interpolated Trotter be composed with LCU or qubitization
to further reduce prefactors, or do the two error-mitigation strategies
interfere?
\emph{Basis.} The paper positions the method as a strict improvement
over single-Trotter but does not compare vs (nor combine with) LCU or
qubitization, which achieve the same Heisenberg limit through a
different mechanism (block-encoding). Whether the two stack
multiplicatively, additively, or destructively is open, and practically
important because qubitization has better constants on structured $H$.
\emph{Next steps.} Implement a hybrid: LCU/qubitization for a coarse
block-encoded evolution, extract expectation values at several
block-encoding accuracies, then Chebyshev-interpolate. Compare gate
counts vs (a) LCU alone at target $\varepsilon$, (b) Chebyshev+Trotter
alone. Classical Qiskit statevector on $\le 6$ qubits.

\item \textbf{Chemistry / lattice benchmarks.}
For chemically or physically relevant Hamiltonians (e.g., 1D Hubbard at
half filling, or a small FeMoco active-space fragment), what is the
actual crossover in gate count where Cheb+$S_2$ beats optimized $S_4$
or higher-order composite ST formulas?
\emph{Basis.} Paper's Sec 5 is on a 2-spin TFIM (trivial commutator
structure, exact diagonalization). For chemistry Hamiltonians with
hundreds of Pauli terms and large $\|[H_1, H_2]\|$, the Cheb cost
prefactor (each node costs $1/s_k$ $S_2$ exponentials) may lose to
higher-order ST despite better asymptotic scaling.
\emph{Next steps.} Instantiate $H$ = 1D Hubbard on $N=4$--$8$ sites
(Jordan-Wigner). Reuse the v2 interpolation loop; compute error vs
total-CNOT-count for Cheb+$S_2$, single $S_4$, single $S_6$, and
Yoshida composites. Report crossover $N$ and gate count.

\item \textbf{Noise robustness.}
How robust is the Chebyshev-interpolated eigenvalue estimate under
realistic noise, i.e.\ when each node estimate carries depolarizing or
amplitude-damping error at hardware-relevant rates
($10^{-3}$--$10^{-4}$ per two-qubit gate)?
\emph{Basis.} The paper is a noiseless-simulator study; extrapolation
of noisy data is generically ill-conditioned (Runge-like amplification
of node error). Whether the exponential-in-$n$ interpolation gain
survives, or whether the noise floor kills it at some crossover $n^*$,
is not addressed --- directly bearing on whether the method is useful
pre-fault-tolerance.
\emph{Next steps.} Wrap the $S_2$ step in Qiskit Aer with a
depolarizing noise model ($p_{2q}\!\in\![10^{-4},10^{-2}]$); compute
expectation values at each $s_k$ node with $M$ shots; fit interpolation
error vs $n$ at multiple $(p_{2q}, M)$. Identify $n^*$ where sampling
plus noise variance overtakes interpolation gain.

\item \textbf{ML-optimized interpolation nodes.}
Can classical ML (small neural surrogate or Gaussian-process
regression) pick data-driven interpolation nodes $s_k$ that outperform
static Chebyshev-1st-kind nodes for a given Hamiltonian family?
\emph{Basis.} Chebyshev nodes are provably optimal for interpolating
smooth functions on an interval in the sup norm, but the true function
$\tilde H_s$ here has extra structure (evenness, dominant Trotter-error
scaling $\sim s^2$) that generic node placement ignores. A learned node
schedule per Hamiltonian family may reduce node count at fixed
precision, especially in the low-$n$ regime ($n=2$--$4$) where each
node is expensive.
\emph{Next steps.} Build a dataset of $\{H_i, s_{\mathrm{grid}},
\mathrm{err}(s)\}$ over 100--1000 random 2--4 qubit local Hamiltonians
(offline Qiskit). Train a small MLP / GP to predict optimal node
placement from $\|[H_1,H_2]\|$ and spectral radius. Compare test-set
interp error vs Chebyshev nodes at fixed $n$. Argo Opus for the ML
training loop; classical Qiskit for data generation.

\end{enumerate}
